\section{Conclusion}
\label{sec:conclusion}

We presented BDM-AOS, a bi-level decomposition matheuristic with adaptive operator selection for the multi-objective Movie Scene Scheduling Problem.
Our work makes four primary contributions:

\begin{enumerate}
    \item A \emph{comprehensive MILP formulation} with 13 constraint families over a three-tier geographic hierarchy (city--town--location), actor--scene--day and location--scene--day linking variables, town-uniqueness and town-continuity logistics constraints, criticality-weighted utilization, and actor consecutive working-day limits---substantially extending prior formulations.

    \item The \emph{Scene Interaction Graph} decomposes scenes into shooting blocks based on actor-sharing, location-sharing, and transfer cost structure, reducing the search space from $n!$ to $K!$ permutations.
    For our largest instance (100 scenes, 24 blocks), this reduction exceeds $10^{134}\times$.

    \item The \emph{bi-level framework} combines NSGA-II over block permutations with a global greedy decoder that enforces all 13 constraint families. The decomposition guides scene ordering while cross-block day sharing preserves cost-competitiveness. This is the first decomposition-based approach for the MSSP.

    \item \emph{Q-learning adaptive operator selection} with two novel problem-specific operators---Location-Aware Crossover (LAX) and Actor-Continuity Mutation (ACM)---adapts the search strategy based on population diversity, convergence rate, and search progress.
\end{enumerate}

Comprehensive experiments on 8 instance sizes (10--100 scenes) with 30 independent seeds per algorithm-instance pair demonstrate the approach's scalability.
A five-variant ablation study quantifies the contribution of each component:
SIG-guided decomposition, adaptive operator selection, and problem-specific operators.
The scalability analysis shows that BDM-AOS's advantages grow with instance size, as the decomposition provides increasingly valuable structure exploitation at larger scales.

\subsection{Future Work}

Several promising directions emerge:

\begin{itemize}
    \item \textbf{Real-world validation:} Apply BDM-AOS to real film scripts (e.g., the 70-scene, 43-actor instance from~\citet{tekin2023}) to validate the hierarchical location model and geographic logistics constraints under real production conditions.

    \item \textbf{Adaptive decomposition:} Re-cluster the SIG periodically during search to exploit evolving population structure, or use hierarchical decomposition (city-level blocks containing town-level sub-blocks) to address the block-granularity limitation.

    \item \textbf{Alternative decoders:} Design decoders that produce genuine cost-makespan trade-offs by allowing variable daily utilization, explicitly trading off town-switching travel days against shooting compactness, and enabling the bi-objective formulation to realize its full potential.

    \item \textbf{Deep RL for AOS:} Replace tabular Q-learning with deep Q-networks for continuous state representation, enabling finer-grained operator selection that generalizes across instance sizes and constraint configurations.

    \item \textbf{Many-objective extension:} Extend to three or more objectives (cost, makespan, actor fatigue, location utilization, town-switching frequency) using NSGA-III~\citep{deb2014}, leveraging the linking variables $v_{a,s,d}$ and $z_{l,s,d}$ for fine-grained objective computation.

    \item \textbf{Uncertainty handling:} Incorporate stochastic elements (weather delays, actor illness, location unavailability) via robust or stochastic programming at the sub-problem level, using the hierarchical structure to model correlated disruptions (e.g., weather affecting all locations within a town).

    \item \textbf{Extended hierarchy:} Model additional tiers (country-level for international productions) and incorporate travel-day constraints between cities, where town continuity (\ref{c12}) generalizes to mandatory rest days between city changes.
\end{itemize}
